\title{Group Sequence Policy Optimization}

\begin{abstract}
We present Group Sequence Policy Optimization (GSPO), a reinforcement learning algorithm designed for training large language models with enhanced stability, efficiency, and performance. In contrast to earlier approaches that utilize token-level importance ratios, GSPO introduces an importance ratio formulated over entire sequences and applies clipping, rewards, and optimization at the sequence level. Our results show that GSPO delivers improved training efficiency and effectiveness relative to the GRPO algorithm, markedly increases stability in Mixture-of-Experts (MoE) RL scenarios, and streamlines RL system design. The adoption of GSPO has played a significant role in the notable advancements seen in the latest Qwen3 models.
\end{abstract}

\section{Introduction}
\label{sec:intro}

Reinforcement learning (RL) has become a foundational approach for advancing the capabilities of language models \citep{o1,dpsk-r1,qwq32b,qwen3}. By employing RL at scale, language models gain proficiency in addressing complex challenges, including advanced mathematics and programming tasks, through extended and more intricate reasoning pathways.

Achieving effective scaling of RL with increased computational resources fundamentally depends on ensuring stable and resilient training dynamics. Yet, leading RL approaches such as GRPO \citep{grpo} frequently encounter pronounced instability during the training of very large language models, with training runs sometimes ending in catastrophic, unrecoverable model failures \citep{qwen3, minimax-m1}. Such instability acts as a major barrier to further advancing the performance limits of language models via prolonged RL optimization.

In this work, we demonstrate that GRPO's instability arises primarily from a flawed use and consequent breakdown of importance sampling weights within its algorithm. This misapplication generates significant variance in training noise, which not only escalates as the response length increases but is also intensified by the application of the clipping mechanism, eventually leading to the collapse of the model.

In response to these fundamental shortcomings, we introduce \textbf{Group Sequence Policy Optimization (GSPO)}, a novel reinforcement learning approach tailored for large language model training. GSPO's primary contribution is a theoretically robust formulation of the importance ratio utilizing sequence likelihood \citep{click}, which is consistent with the foundational concept of importance sampling. Furthermore, GSPO calculates normalized rewards to serve as advantages across several responses to a given query, thereby preserving the correspondence between rewards assigned at the sequence level and the optimization process.

Our experimental results indicate that GSPO substantially outperforms GRPO with respect to training stability, efficiency, and overall effectiveness. Notably, GSPO intrinsically addresses stability issues commonly encountered during the reinforcement learning training of large Mixture-of-Experts (MoE) models, thereby obviating the requirement for intricate stabilization techniques and suggesting a pathway toward more streamlined RL frameworks. These advantages of GSPO were key drivers behind the remarkable performance enhancements observed in the most recent Qwen3 models. We anticipate that GSPO will serve as a reliable and extensible algorithmic platform, fostering the progression of large-scale RL training for language models.

\section{Preliminaries}

\paragraph{Notation}
Throughout this work, we refer to an autoregressive language model with parameters $\theta$ as a policy $\pi_\theta$. We represent a query by $x$, and the set of all queries is denoted by $\mathcal{D}$. For a given query $x$ and a generated response $y$, the probability assigned to $y$ by policy $\pi_\theta$ is expressed as $\pi_\theta (y | x)=\prod_{t=1}^{|y|} \pi_\theta (y_t | x, y_{<t} )$, with $|y|$ indicating the length of $y$ in tokens. The quality or suitability of a query-response pair $(x, y)$ is assessed via a verifier $r$, which provides a reward score $r(x, y) \in [ 0, 1 ]$.

\paragraph{Proximal Policy Optimization (PPO)} Given data collected under the previous policy $\pi_{\theta_\text{old}}$, PPO \citep{ppo} restricts policy changes by enforcing updates to remain close to the preceding policy, utilizing a clipping strategy. In particular, PPO optimizes the policy using the following loss function (excluding the KL regularization term for simplicity, since it is not central to the current work):

\begin{align}
\mathcal{J}_\text{PPO}(\theta) = \mathbb{E}_{ x \sim \mathcal{D},\, y \sim \pi_{\theta_\text{old}}( \cdot | x) }
\left[ \frac{1}{|y|} \sum_{t=1}^{|y|} 
\min \left( w_{t}(\theta) \widehat{A}_{t},  \, \mathrm{clip} \left( w_{t}(\theta), 1 - {\varepsilon}, 1 + {\varepsilon}\right) \widehat{A}_{t} \right)
\right],
\end{align}

Here, the token $y_t$'s importance ratio is given by $w_{t}(\theta) = \frac{ \pi_{\theta} (y_{t} | x, y_{<t}) }{ \pi_{\theta_\text{old}} (y_{t} | x,y_{<t})}$; the advantage $\widehat{A}_{t}$ corresponding to $y_t$ is approximated using a separate value model, and $\varepsilon$ denotes the interval within which the importance ratio is clipped.

A primary obstacle encountered in the practical application of PPO is its significant dependence on the value model. This model often matches the policy model in scale, resulting in substantial demands on both memory and computational resources. Additionally, the performance of PPO is closely tied to the accuracy of its value predictions. Achieving dependable value estimation is intrinsically difficult, and the difficulties are amplified when seeking models capable of handling longer outputs and increasingly sophisticated tasks.

\paragraph{Group Relative Policy Optimization (GRPO)} As described in \citep{grpo}, GRPO eliminates the reliance on a value model by evaluating the advantage of a given response relative to others in a set that address the identical query. In particular, GRPO seeks to maximize the following objective:

\begin{align}
\label{equ:grpo}
\mathcal{J}_\text{GRPO}(\theta) = \mathbb{E}_{ x \sim \mathcal{D},\, \{y_i\}_{i=1}^G \sim \pi_{\theta_\text{old}}( \cdot | x) }
\left[ \frac{1}{G} \sum_{i=1}^{G} \frac{1}{|y_i|} \sum_{t=1}^{|y_i|} 
\min \left( w_{i,t}(\theta) \widehat{A}_{i,t},  \, \mathrm{clip} \left( w_{i,t}(\theta), 1 - {\varepsilon}, 1 + {\varepsilon}\right) \widehat{A}_{i,t} \right)
\right],
\end{align}

Here, $G$ refers to the total number of responses produced for each query $x$ (that is, the group size). The importance ratio $w_{i,t}(\theta)$ and the advantage $\widehat{A}_{i,t}$ associated with token $y_{i,t}$ are defined as follows:

\begin{align}
    w_{i,t}(\theta)=\frac{ \pi_{\theta} (y_{i,t} | x, y_{i,<t}) }{ \pi_{\theta_\text{old}} (y_{i,t} | x,y_{i,<t})},\quad\ 
    \widehat{A}_{i,t} = \widehat{A}_{i} = \frac{r(x, y_i) - \mathrm{mean} \left( \{ r(x, y_i) \}_{i=1}^G \right) }{ \mathrm{std} \left( \{ r(x, y_i) \}_{i=1}^G \right) },
\end{align}

respectively, with each token in $y_i$ possessing the identical advantage value $\widehat{A}_{i}$.

\section{Motivation}
\label{sec:motivation}

As model sizes increase, along with factors such as sparsity (as in Mixture-of-Experts architectures) and longer output sequences, there is a demand for large rollout batch sizes to fully leverage computational hardware in RL settings. In order to make learning more sample-efficient, a widely adopted approach is to divide a sizable rollout batch into several mini-batches for performing gradient updates. However, this approach creates an off-policy learning scenario: the generated responses $y$ are drawn from a previous policy $\pi_{\theta_\text{old}}$, as opposed to the current policy $\pi_\theta$ that is actively being updated. This condition underscores the importance of employing clipping techniques in PPO and GRPO, which serve to filter out samples that are excessively off-policy and thus mitigate their influence on the gradient computation.

Although strategies such as clipping are designed to address off-policy discrepancies, our analysis reveals a deeper, foundational flaw within GRPO: \textit{the objective itself is not well-defined}. This challenge is especially pronounced during the training of large-scale models on tasks requiring extended responses, where it can result in severe degradation and eventual collapse of model performance. The root of this ill-posedness in the GRPO objective can be traced to incorrect usage of importance sampling weights. Importance sampling, in principle, aims to compute the expectation of a function $f$ under a desired distribution $\pi_\text{tar}$ by appropriately reweighting data sampled from an alternative, behavior distribution $\pi_\text{beh}$:

\begin{align}
\mathbb{E}_{z \sim \pi_\text{tar}} \left[ f(z) \right] = \mathbb{E}_{z \sim \pi_\text{beh}} \left[ \frac{ \pi_\text{tar} (z) }{ \pi_\text{beh} (z)} \, f(z) \right].
\end{align}

Importantly, the effectiveness of correcting for distributional discrepancies using the importance weight $\frac{ \pi_\text{tar} (z) }{ \pi_\text{beh} (z)}$ hinges on averaging across a large number of samples ($N \gg 1$) drawn from the behavior distribution $\pi_\text{beh}$.

By contrast, GRPO utilizes the importance weighting factor $\frac{ \pi_{\theta} (y_{i,t} | x, y_{i,<t}) }{ \pi_{\theta_\text{old}} (y_{i,t} | x,y_{i,<t})}$ at each individual token step $t$. This approach relies solely on a single token sample $y_{i,t}$ from the next-token distribution $\pi_{\theta_\text{old}}( \cdot | x, y_{i, <t})$, which undermines its ability to effectively correct for distribution shifts. Rather than fulfilling this corrective function, the method injects substantial high-variance noise into the gradient updates during training, and this issue intensifies across lengthy sequences, particularly due to the influence of the clipping procedure. Empirically, we have found that such conditions may precipitate a model collapse, which tends to be permanent. Attempts to recover training progress after such a collapse—including reverting to earlier checkpoints, extensive hyperparameter (such as clipping interval) adjustments, increasing the generated sequence length, or altering the RL query mechanism—have not succeeded.

The preceding analysis highlights a key limitation inherent in GRPO's structure. The ineffectiveness of importance weighting at the token level reveals a crucial insight: \textbf{the granularity of the optimization target must correspond to that of the reward signal}. Because rewards are distributed over whole sequences rather than individual tokens, implementing off-policy adjustments at the token level proves inadequate. Consequently, we are compelled to abandon token-level objectives in favor of leveraging importance weights and conducting optimization at the \textit{sequence level}.

\section{Algorithm}

\subsection{GSPO: Group Sequence Policy Optimization}

Although the use of token-level importance weights $\frac{ \pi_{\theta} (y_{i,t} | x, y_{i,<t}) }{ \pi_{\theta_\text{old}} (y_{i,t} | x,y_{i,<t})}$ presents difficulties within GRPO, our findings indicate that, for language generation tasks, the \textit{sequence-level} importance weight $\frac{ \pi_{\theta} (y | x) }{ \pi_{\theta_\text{old}} (y | x)}$ holds a well-defined theoretical significance. Specifically, this ratio quantifies the extent to which a response $y$—sampled from $\pi_{\theta_\text{old}} (\cdot | x)$—differs from what would be produced under $\pi_{\theta} (\cdot | x)$, thus directly corresponding to the sequence-level reward. Furthermore, this measure effectively captures the notion behind the clipping procedure.

Building upon this simple insight, we introduce the \textbf{Group Sequence Policy Optimization (GSPO)} algorithm. GSPO optimizes a sequence-level objective defined as follows:

\begin{align}
\mathcal{J}_\text{GSPO} (\theta) =
\mathbb{E}_{ x \sim \mathcal{D},\, \{y_i\}_{i=1}^G \sim \pi_{\theta_\text{old}}( \cdot | x) }
\left[ 
\frac{1}{G} \sum_{i=1}^{G}
\min \left( s_{i}(\theta)  \widehat{A}_{i},  \, \mathrm{clip} \left( s_{i}(\theta), 1 - {\varepsilon}, 1 + {\varepsilon} \right) \widehat{A}_{i} \right) 
\right],
\label{equ:gspo}
\end{align}

in which we utilize advantage estimation methods that are structured around groups:

\begin{align}
\widehat{A}_{i} = \frac{r(x, y_i) - \mathrm{mean} \left( \{ r(x, y_i) \}_{i=1}^G \right) }{ \mathrm{std} \left( \{ r(x, y_i) \}_{i=1}^G \right) },
\end{align}

and establish the importance ratio $s_{i}(\theta)$ utilizing sequence likelihood as described in \citet{click}:

\begin{align}
s_{i}(\theta) = \left( \frac{ \pi_{\theta} (y_i | x) }{ \pi_{\theta_\text{old}} (y_i | x)} \right)^{\frac{1}{|y_i|}}
=
\exp \left( \frac{1}{|y_i|} \sum_{t=1}^{|y_i|} \log \frac{ \pi_{\theta} (y_{i,t} | x, y_{i,<t}) }{ \pi_{\theta_\text{old}} (y_{i,t} | x,y_{i,<t})} \right).
\end{align}

Consequently, GSPO implements clipping at the response level rather than at the token level, thereby filtering out highly ``off-policy'' samples during gradient computation, which is more consistent with both sequence-level reward assignment and the corresponding optimization. Additionally, we utilize length normalization in $s_{i}(\theta)$ to minimize variance and to constrain $s_{i}(\theta)$ within a consistent numerical interval. Without such normalization, small likelihood shifts in individual tokens may induce significant volatility in the sequence-level importance ratio, necessitating different clipping thresholds for responses of varying lengths. Furthermore, we observe that the clipping intervals used in GSPO generally differ, often substantially, from those used in prior methods such as GRPO, as a result of the inherently different definitions of importance ratios in these approaches.

\subsection{Gradient Analysis}
\label{subsec:gradient}

The gradient of the GSPO objective can be computed in the following manner (note that we exclude clipping for simplicity):

\begin{align}
\nabla_{\theta} \mathcal{J}_\text{GSPO} (\theta)
=&\ 
\nabla_{\theta} \mathbb{E}_{ x \sim \mathcal{D},\, \{y_i\}_{i=1}^G \sim \pi_{\theta_\text{old}}( \cdot | x) }
\left[ 
\frac{1}{G} \sum_{i=1}^{G} s_{i}(\theta) \widehat{A}_{i}
\right] \\
= &\ 
\mathbb{E}_{ x \sim \mathcal{D},\, \{y_i\}_{i=1}^G \sim \pi_{\theta_\text{old}}( \cdot | x) }
\left[ 
\frac{1}{G} \sum_{i=1}^{G}
s_{i}(\theta) \widehat{A}_{i}
\cdot \nabla_{\theta} \log s_{i}(\theta)
\right] \\
=&\ 
\mathbb{E}_{ x \sim \mathcal{D},\, \{y_i\}_{i=1}^G \sim \pi_{\theta_\text{old}}( \cdot | x) }
\left[ 
\frac{1}{G} \sum_{i=1}^{G} 
\left( \frac{ \pi_{\theta} (y_i | x) }{ \pi_{\theta_\text{old}} (y_i | x)} \right)^{\frac{1}{|y_i|}} \widehat{A}_{i} 
\cdot \frac{1}{|y_i|} \sum_{t=1}^{|y_i|} \nabla_{\theta} \log \pi_{\theta} (y_{i,t} | x, y_{i,<t}) 
\right].
\label{equ:gspo_grad}
\end{align}

To facilitate comparison, the gradient of the GRPO objective is given below (recall that $\widehat{A}_{i,t} = \widehat{A}_{i}$):

\begin{align}
\nabla_{\theta} \mathcal{J}_\text{GRPO} (\theta) 
=&\ 
\nabla_{\theta} \mathbb{E}_{ x \sim \mathcal{D},\, \{y_i\}_{i=1}^G \sim \pi_{\theta_\text{old}}( \cdot | x) }
\left[ \frac{1}{G} \sum_{i=1}^{G} \frac{1}{|y_i|} \sum_{t=1}^{|y_i|} 
w_{i,t}(\theta) \widehat{A}_{i,t}
\right] \\
=&\
\mathbb{E}_{ x \sim \mathcal{D},\, \{y_i\}_{i=1}^G \sim \pi_{\theta_\text{old}}( \cdot | x) }
\left[ \frac{1}{G} \sum_{i=1}^{G} \widehat{A}_{i} 
\cdot \frac{1}{|y_i|} \sum_{t=1}^{|y_i|} 
\frac{ \pi_{\theta} (y_{i,t} | x, y_{i,<t}) }{ \pi_{\theta_\text{old}} (y_{i,t} | x,y_{i,<t})} 
\nabla_{\theta} \log \pi_{\theta} (y_{i,t} | x, y_{i,<t})  
\right].
\end{align}

The principal difference between GSPO and GRPO is rooted in \textit{the approach each method uses to assign weights to the gradients of token log likelihoods}. GRPO utilizes an ``importance weight'' $\frac{ \pi_{\theta} (y_{i,t} | x, y_{i,<t}) }{ \pi_{\theta_\text{old}} (y_{i,t} | x,y_{i,<t})}$ for each token, meaning that tokens contribute unequally based on these values. These weights, which may fall within the ranges $(0, 1+\varepsilon]$ when $\widehat{A}_{i} >0$, or $[1-\varepsilon, +\infty)$ when $\widehat{A}_{i} < 0$, are substantial and not to be ignored; as training advances, the aggregation of such disparities can introduce instability and unforeseen effects. GSPO, on the other hand, assigns uniform weight to every token within a response, thereby removing the particular instability introduced by the non-uniform weighting characteristic of GRPO.

\subsection{GSPO-token: A Token-level Objective Variant}

In contexts such as multi-turn RL, a more detailed advantage modification than at the sequence level may be necessary. Therefore, we propose a token-level extension of GSPO, referred to as \textbf{GSPO-token}, which facilitates the customization of advantages for individual tokens:

\begin{align}
\mathcal{J}_\text{GSPO-token}(\theta)
=
\mathbb{E}_{ x \sim \mathcal{D},\, \{y_i\}_{i=1}^G \sim \pi_{\theta_\text{old}}( \cdot | x) }
\left[ 
\frac{1}{G} \sum_{i=1}^{G} \frac{1}{|y_i|} \sum_{t=1}^{|y_i|} 
\min \left( s_{i,t}(\theta) \widehat{A}_{i,t},  \, \mathrm{clip} \left( s_{i,t}(\theta), 1 - {\varepsilon}, 1 + {\varepsilon} \right) \widehat{A}_{i,t} \right) 
\right],
\label{equ:gspo-token}
\end{align}

in which

\begin{align}
s_{i,t}(\theta) 
= \mathrm{sg} \left[ s_{i}(\theta) \right]  \cdot \frac{ \pi_{\theta} (y_{i,t} | x, y_{i,<t}) }{ \mathrm{sg} \left[ \pi_{\theta} (y_{i,t} | x, y_{i,<t}) \right] },
\end{align}

and $\mathrm{sg}[\cdot]$ indicates extracting the numerical value while preventing gradient flow, equivalent to applying the \texttt{detach} function in PyTorch. The derivative of GSPO-token is given by:

\begin{align}
\nabla_{\theta} \mathcal{J}_\text{GSPO-token}(\theta)
=&\ 
\nabla_{\theta} \mathbb{E}_{ x \sim \mathcal{D},\, \{y_i\}_{i=1}^G \sim \pi_{\theta_\text{old}}( \cdot | x) }
\left[ 
\frac{1}{G} \sum_{i=1}^{G}
\frac{1}{|y_i|} \sum_{t=1}^{|y_i|} 
s_{i,t}(\theta) \widehat{A}_{i,t}
\right] \\
=&\ 
\mathbb{E}_{ x \sim \mathcal{D},\, \{y_i\}_{i=1}^G \sim \pi_{\theta_\text{old}}( \cdot | x) }
\left[ 
\frac{1}{G} \sum_{i=1}^{G} s_i (\theta) \cdot \frac{1}{|y_i|} \sum_{t=1}^{|y_i|} 
\widehat{A}_{i,t} \frac{ \nabla_{\theta} \pi_{\theta} (y_{i,t} | x, y_{i,<t}) }{ \pi_{\theta} (y_{i,t} | x, y_{i,<t}) }
\right] \\
=&\ 
\mathbb{E}_{ x \sim \mathcal{D},\, \{y_i\}_{i=1}^G \sim \pi_{\theta_\text{old}}( \cdot | x) }
\left[ 
\frac{1}{G} \sum_{i=1}^{G}  \left( \frac{ \pi_{\theta} (y_i | x) }{ \pi_{\theta_\text{old}} (y_i | x)} \right)^{\frac{1}{|y_i|}}
\cdot \frac{1}{|y_i|} \sum_{t=1}^{|y_i|}
\widehat{A}_{i,t} \nabla_{\theta} \log \pi_{\theta} (y_{i,t} | x, y_{i,<t})  
\right].
\label{equ:gspo-token_grad}
\end{align}

Observe that $\frac{ \pi_{\theta} (y_{i,t} | x, y_{i,<t}) }{ \mathrm{sg} \left[ \pi_{\theta} (y_{i,t} | x, y_{i,<t}) \right] }$ evaluates to 1, which means that, in terms of their actual values, $s_{i,t}(\theta)$ and $s_{i}(\theta)$ are equivalent. By examining Equation~\eqref{equ:gspo} alongside~\eqref{equ:gspo-token}, as well as Equation~\eqref{equ:gspo_grad} in comparison with~\eqref{equ:gspo-token_grad}, it becomes evident that GSPO and GSPO-token yield the same numerical results for the optimization objective, clipping criterion, and formal gradient, provided that all token advantages within the response $y_i$ are set equally (that is, $\widehat{A}_{i,t} = \widehat{A}_{i}$). Nevertheless, GSPO-token offers additional versatility by permitting the allocation of different advantage values to each token individually.

\section{Experiments and Discussion}

\subsection{Empirical Results}

We utilize a cold-start model fine-tuned from Qwen3-30B-A3B-Base and present both the training reward trajectories and the evaluation metrics for the model across the AIME'24 benchmark (mean Pass@1 with 32 samples), LiveCodeBench (202410-202502, mean Pass@1 with 8 samples), and CodeForces (Elo Rating). In reinforcement learning (RL) training, each rollout batch is subdivided into four mini-batches for the purposes of gradient updating. For GSPO, the clipping boundaries in Equation~\eqref{equ:gspo} are assigned values of 3e-4 (left) and 4e-4 (right). Our baseline comparator is GRPO, where the clipping bounds in Equation~\eqref{equ:grpo} are established at 0.2 (left) and 0.27 (right), parameters that we have meticulously optimized to provide a rigorous comparative evaluation. It is important to highlight that GRPO relies on the Routing Replay method to achieve reliable convergence in MoE RL settings, as will be discussed further in \S~\ref{subsec:routing_replay}, whereas \textit{GSPO eliminates the dependency on this approach}.

As depicted in Figure~\ref{fig:results}, the GSPO training process remains stable across the duration of training. Notably, \textbf{GSPO consistently enhances performance by scaling up training compute, periodically refreshing the query set, and increasing generation length}. Additionally, GSPO surpasses GRPO in terms of training efficiency, attaining higher training accuracy and improved benchmark outcomes given equivalent training compute and query consumption. Lastly, we have effectively leveraged GSPO for RL-based training of the newest Qwen3 models, providing strong evidence for GSPO's ability to unlock the scaling potential of reinforcement learning for large language models.

\subsection{Curious Observation on Clipping Fractions}

An important differentiation between GSPO and GRPO lies in GSPO's strategy of clipping entire sequences, rather than focusing on individual tokens. As depicted in Figure~\ref{fig:clipping}, this approach results in a disparity of roughly two orders of magnitude in the proportion of tokens that are clipped between GSPO and GRPO; notably, modifying the clipping thresholds does not bridge this substantial gap. Nevertheless, even though GSPO discards a much greater number of tokens—leading to fewer data points for gradient calculation or training—it consistently surpasses GRPO in training efficiency. This seemingly paradoxical outcome, where discarding a larger proportion of data enhances training efficiency, underscores that GRPO’s gradient estimation at the token level is fundamentally noisy and suboptimal for leveraging sampled data. By comparison, GSPO's sequence-level mechanism generates a more stable and productive learning signal.

\subsection{Benefit of GSPO for MoE Training}
\label{subsec:routing_replay}

\paragraph{Background}
The sparse activation characteristic inherent to MoE models, as opposed to the dense structure of conventional models, gives rise to specific stability issues during RL training. Notably, our observations indicate that, when employing the GRPO algorithm, the phenomenon of \textit{expert-activation volatility} can obstruct the stable convergence of RL training in MoE architectures. Specifically, following one or more gradient steps, the subset of experts engaged in producing the same model output may shift considerably. For instance, in the case of the 48-layer Qwen3-30B-A3B-Base model, approximately 10\% of the experts chosen by the updated policy $\pi_{\theta}$ for a particular rollout diverge from those selected by the previous policy $\pi_{\theta_\text{old}}$. This volatility is particularly accentuated in MoE models with increased depth and results in severe oscillations of the token-level importance weights $w_{i,t}(\theta) = \frac{ \pi_{\theta} (y_{i,t} | x, y_{i,<t}) }{ \pi_{\theta_\text{old}} (y_{i,t} | x,y_{i,<t})}$, thereby rendering these weights unreliable, as examined in \S~\ref{sec:motivation} and~\ref{subsec:gradient}. As a result, the standard convergence behavior of RL training is adversely impacted.

\paragraph{Our Previous Approach} In our earlier work addressing this issue, we adopted the \textbf{Routing Replay} training methodology. This method involves storing the selected experts activated under $\pi_{\theta_\text{old}}$, then using these same routing decisions, or ``replaying'' them, under the current policy $\pi_{\theta}$ during the calculation of importance ratios, defined as $w_{i,t}(\theta) = \frac{ \pi_{\theta} (y_{i,t} | x, y_{i,<t}) }{ \pi_{\theta_\text{old}} (y_{i,t} | x,y_{i,<t})}$. Through this approach, both $\pi_{\theta} (y_{i,t} | x, y_{i,<t})$ and $\pi_{\theta_\text{old}} (y_{i,t} | x,y_{i,<t})$ for a given token $y_{i,t}$ utilize an identical set of experts, which stabilizes the resulting token-level importance ratios. This in turn preserves consistent optimization of the chosen subnetwork throughout successive gradient steps. As illustrated in Figure~\ref{fig:routing_replay}, Routing Replay plays a crucial role in achieving proper convergence behavior during the GRPO training procedure for MoE models.

\paragraph{Benefit of GSPO} While Routing Replay allows GRPO training of MoE models to achieve stable convergence, it does so at the cost of increased memory usage and communication demands, in addition to potentially constraining the effective capacity of the MoE architecture. On the other hand, as depicted in Figure~\ref{fig:results}, GSPO removes the reliance on Routing Replay, seamlessly computes the importance ratios $s_i(\theta)$ in the usual manner, and maintains robust and stable convergence throughout optimization. The central idea is that GSPO evaluates only the sequence-level likelihood (i.e., $\pi_{\theta} (y_i | x)$) and does not depend on the likelihood of individual tokens (i.e., $\pi_{\theta} (y_{i,t} | x, y_{i,<t})$). Because the MoE model's language modeling capabilities are preserved at all times, the sequence likelihood remains consistently stable. Ultimately, GSPO directly addresses the issue of instability in expert activation for MoE models, thereby removing the need for workaround strategies such as Routing Replay. As a result, the training process becomes both more straightforward and reliable, and the model can make full use of its underlying capacity without artificial limitations.

\subsection{Benefit of GSPO for RL Infrastructure}

Due to differences in numerical precision between training engines such as Megatron and inference engines like SGLang and vLLM, it is common practice to recompute the likelihoods of sampled outputs under the previous policy $\pi_{\theta_\text{old}}$ using the training engine. In contrast, GSPO performs optimization based solely on sequence-level likelihoods rather than at the token-level. Notably, sequence-level likelihoods exhibit a higher tolerance to precision mismatches. As a result, GSPO enables the direct utilization of likelihoods produced by the inference engine during optimization, removing the requirement for additional recomputation by the training engine. This approach proves particularly advantageous in cases involving partial rollouts, multi-turn reinforcement learning, or training-inference disaggregated architectures.

\section{Conclusion}

We introduce Group Sequence Policy Optimization (GSPO), an innovative reinforcement learning algorithm designed for training large language models. GSPO leverages the core concept of importance sampling by constructing importance ratios grounded in sequence likelihood, enabling sequence-level clipping, reward assignment, and policy optimization. Compared to GRPO, GSPO offers marked improvements in training stability, efficiency, and effectiveness, particularly excelling in large-scale reinforcement learning for MoE models. This approach forms the basis for the notable advancements seen in recent Qwen3 models. As a scalable framework, GSPO enables further expansion of RL methods, positioning us to pursue significant progress in artificial intelligence.

\end{document}